\section{Methodology}\label{sec:method}

\subsection{Real-engine harness}\label{sec:harness}
All real-engine runs used one NVIDIA A100 80GB PCIe GPU per job on the Nova SLURM cluster at Iowa State University, with SGLang 0.5.22 (FlashInfer attention, chunked prefill of 8192 tokens, at most 256 running requests, LRU eviction) serving Qwen2.5-7B-Instruct~\citep{qwen2024qwen25}. The KV pool was capped at 300,000 tokens or, in the native-pool runs, sized by SGLang at 955,757 tokens. Decoding was greedy, 100 tokens per turn (Appendix~\ref{app:protocol} lists every setting). The vLLM 0.32.0 runs used the same model, GPU, client and workload with a 300,000-token pool (\S\ref{sec:vllm}).

An open-loop client on the same node sends each turn as a completion request whose prompt is a list of token ids: random content from a seeded generator, extended each turn by the token ids the server returned plus 550 new tool tokens, so a hit covers exactly the previous prompt and output; during a tool call the client sleeps. Cached tokens come from the server's usage report, which the engine's counters and log confirmed in the closed-population runs (Stage~C, \S\ref{sec:stages}). Before every run the cache is flushed and checked with a canary: a 1024-token request that reused 1023 tokens before the flush (1008 on vLLM, which caches whole 16-token blocks) must reuse at most 15 after; it reused none. In the cache-off E0a runs the check only confirms that the engine is idle. Each arm of a registered test ran in its own job on a fresh server.

The BURST arm adds 200 fresh sessions within 60\,s at $t_b=1200$\,s to the Poisson stream; NOBURST is the same stream without them; CAPPED is BURST plus a client-side FIFO gate admitting at most $\ncap$ sessions, each holding its slot until its last turn ends (the gate of Proposition~\ref{prop:cap}). The burst equals 12 to 38 minutes of base arrivals at the registered rates, and its 2000-token first prompts alone (400k tokens) exceed the 300k-token pool. We took its size from our earlier simulator, where bursts of 100 or more trapped every work-conserving policy, so the test asks whether a burst large enough to cross the basin boundary traps the engine, not how likely such a burst is (in the calibrated simulator 25 sessions never trap and 50 trap about half the time, \S\ref{sec:robust}). The arms of one seed share arrival times, turn counts and tool times (common random numbers), so BURST differs from NOBURST by the burst and CAPPED from BURST by the gate, up to GPU timing nondeterminism: each arm ran on a fresh server in its own job, and the arms of one seed did not always share a node. Seeds 0, 1 and 2 are also reused across rates, pools, policy cells and engines. For a given seed the 200 burst sessions are identical in every configuration, the arrival times at a lower rate are a subset of those at a higher rate, and the $i$-th session to arrive has the same turn count, tool times and client-generated tokens everywhere. Runs that share a seed are therefore not independent across rates, pools or engines, and the BURST arm of the registered vLLM test repeats the workloads of the exploratory vLLM runs.

\subsection{Staged protocol}\label{sec:stages}
Each stage's inputs came from the previous one (job IDs in Table~\ref{tab:stages}). \emph{Stage~0}, a smoke test of server and client on one GPU, came first; its server logs and those of Stage~C calibrated the simulator (\S\ref{sec:sim}). \emph{Stage~C}, the closed-population sweep, measured $\mu(n)$ with all $n$ sessions starting at $t=0$, one run per $n$, a 900\,s warmup and an 1800\,s window. \emph{E0a}, the all-miss calibration, measured $\lmiss$ with the radix cache disabled; the reference rate $\lref$ is its turns-based value, 0.0715\sps, and arrival rates are multiples of it. Each engine and pool has its own $\lref$ (0.0718\sps{} with the native pool, 0.0657 on vLLM), and numeric multiples of $\lmiss$ in this paper are multiples of that pool's $\lref$. \emph{E0b}, the from-empty sweep, ran Poisson arrivals from an empty engine for 3000\,s per rate, with arrival seed 0 at every rate; a run is stable if its final queue is below 50 and its queue rises by at most 10 over the last 1500\,s. The registration was then filled and committed, and \emph{E2}, the registered paired test, ran the three arms. Afterwards we repeated the sequence with the native pool (Stage~C, E0a, E0b, a second registration and \emph{native E2}) and ran an exploratory policy grid. On vLLM we first measured $\lmiss$ and ran the BURST arm alone under an exploratory registration, and later ran its own Stage~C and E0b, a third registration and its paired test (\S\ref{sec:vllm}).

\subsection{Decision rule and registration}\label{sec:prereg}
Each run of a registered test lasts 8700\,s and is judged over $[t_b+4500,\,8700]$\,s; in the rule, $\lmiss$ is the cell's $\lref$. A BURST run is TRAPPED if $\lambda>\lmiss$ and (a) its uncached prefill token rate is within $\pm15\%$ of the cell's E0a all-miss rate (11,243 tokens/s at 300k) and (b) the slope of its backlog exceeds $0.5(\lambda-\lmiss)$, the slope being the least-squares slope of 10\,s bin means of the sessions in the system (arrived and not ended, including any held at the gate) over the window. A run is RECOVERED if over the last 300\,s its hit fraction exceeds 0.8 and its outstanding backlog (requests in flight plus sessions held at the gate) is below 20; NOBURST is good if RECOVERED. A run that is neither is DRAINING if its slope is negative and INDETERMINATE otherwise; the labels take precedence in the order trapped, recovered, draining, indeterminate. A rate PASSES if, in at least two of three seeds, BURST is trapped, NOBURST good and CAPPED recovered in that same seed. Criterion (a) uses tokens because a trapped engine with a growing backlog serves mostly young sessions, whose turns are short, so at the all-miss token rate it completes more turns per second than the stationary all-miss mix, and sessions faster than $\lmiss$ (\S\ref{sec:fluid}): on traps simulated in our earlier simulator (\S\ref{sec:sim}), a session-rate test fired on at most 1 of 29 at 2 $\lmiss$ and none of 39 at 2.5, while the token test fired on all 84 traps at 1.5--2.5 $\lmiss$.

The registration was filled by formula from committed copies of the Stage~C curve, the E0a calibration and the E0b bracket. The window was $(1,\,2.0)\,\lref$, its upper edge E0b's largest all-stable multiple (below $\mu^*/\lref=3.26$). The rates followed a rule proposed in the runbook and adopted in the project status notes on 2026-10-22, the day before E0b was submitted (Appendix~\ref{app:prereg}): one third and two thirds of the way across the window, rounded to 0.05, which gives 1.35 and 1.65 $\lref$ (0.0966 and 0.1180\sps). The cap was the smallest $n$ with $\mu(n)\ge0.95\mu^*$: $\ncap=20$, since $\mu(20)=0.224$ is about 4\% below $\mu(24)=0.233$, inside the rule's 5\%. The file was committed and reached the project's GitHub repository at 15:09:54 CDT on 2026-10-24, before the E2 jobs were submitted at 15:13 (Table~\ref{tab:timeline}). We call a test pre-registered when its registration file was committed and pushed to GitHub before its runs; this is not an entry in a third-party registry. Commit times are set by the committer, so the external record is GitHub's push log, archived as \fpath{provenance/github_activity_2026-10.json} (Appendix~\ref{app:timeline}). The job script refuses to start unless the file is committed and unmodified and the job's cell, rates, cap, reference rates, hardware and protocol match it, and writes its SHA-256 into every run; the analysis calls a verdict decisive only for the registered cell and rates with every run bound to the file and started after its commit (Appendix~\ref{app:prereg}); the label records this binding, not statistical strength, and the analysis checks commit order, not push time, so its own label ``pre-registered'' does not by itself meet our definition. The native registration was filled and enforced the same way and reached GitHub 10 minutes before its jobs were submitted (\S\ref{sec:e2native}). The vLLM registration was filled from vLLM's own Stage~C curve, E0a calibration and E0b bracket, and does not meet the definition: it was committed 37\,s before its 18 jobs were submitted, and the cluster's record of the submission names the file but not its hash, but it reached GitHub 2\,min after the submission, when two of the jobs had started and none could have ended, so we call this test registered, not pre-registered (\S\ref{sec:vllm}).

\subsection{The calibrated simulator}\label{sec:sim}
Our discrete-event simulator models a continuous-batching engine: iterations of at most 8192 prefill tokens and 256 sequences; admission reserves a request's end-of-turn KV; idle KV of sessions in a tool call or waiting is evicted whole, least recently used first; waiting requests pin nothing; a finished session's KV is freed at once, an oracle, since a real engine keeps it as evictable cache. An iteration costs a fixed time plus a per-token term plus attention terms. The five cost parameters were fitted only to SGLang's per-batch timing lines in the Stage~0 and Stage~C server logs (decode $R^2=0.992$, median error 0.75\%; prefill $R^2=0.998$, 1.3\%), never to $\mu(n)$. It runs mixed prefill--decode iterations or, like SGLang, separate iterations with prefill priority; the primary (mixed) mode was the one closer to the real Stage~C $\mu(n)$ under a pre-declared metric (mean relative error 0.104 against 0.127), a margin carried by one cell, and both modes' predictions are reported. It does not model partial prefixes, the waiting-prefix refresh or the overlap scheduler; reserving whole turns at admission is an upper bound on what paged engines hold (in round 5 of the earlier simulator, on-demand allocation with preemption and SGLang-like ratio reservation kept the curve and the trap).

\paragraph{Validation.} The declared level test, simulated $\mu$ within 15\% of the real value in each of 8 Stage~C cells, failed in both modes: mixed batching missed one cell, FCFS $n=48$, by $+42\%$ (0.070 against 0.049\sps), separate batching two. The shape was reproduced: the hit fraction crosses 0.5 at $n=32.5$ against 32.2, the plateau ($n=16$--24) agrees within 9.2\% and $n=300$ within 2\%, and $\mu(300)/\mu^*$ is 0.50 against 0.46. The single real run at $n=48$ dipped below the simulator's all-miss level; its turns-based estimate, 0.063\sps, is within 15\%. Keeping a finished session's KV as evictable cache, a labelled sensitivity arm, left mixed batching the closer mode (0.118 against 0.146) but lowered the plateau ($-15.5\%$ at $n=24$; 5 of 8 cells within 15\%). In the earlier simulator the same change moved the orderings and pins of \S\ref{sec:mech} by at most 0.025 $\lmiss$ but stalled the CONCUR surrogate at its minimum cap in 12 of 18 cells; the static caps were not rerun.

\paragraph{Blind predictions.} The calibrated simulator then predicted, in both modes (mixed primary), B1: $\lmiss$ and the all-miss token rate; B2: the fraction of from-empty runs that stay stable at each rate; B3: the E2 outcome at each rate from 0.09 to 0.16\sps; B4: the closed curve with 600k- and 900k-token pools. They were made by LLM coding agents (Claude Sonnet sessions, each in its own copy of the repository at a commit without the E0a or E0b results), which were instructed not to contact the cluster and not to open the E0a or E0b result directories or the primary copy of the repository holding all results; their only real data were the Stage~0 and Stage~C runs on which the simulator was calibrated and validated. They were not technically barred: their tool-call transcripts, which are not in the repository, were checked afterwards by the orchestrating Claude Code session, which found no cluster command and no access to E0b before the predictions reached GitHub. They reached GitHub at 14:53 CDT on 2026-10-24, before any E2 job was submitted and before any native-pool run, so B3 and B4 are blind in time. E0b had already finished on the cluster, and its results were pulled 13 minutes later, so B2 is blind only by instruction and commit order (Table~\ref{tab:timeline}). B1 is not blind in time (E0a had been committed), and it is close to the calibration data: the cost model was fitted to Stage~C per-batch prefill timing, and the all-miss capacity is essentially a prefill-token rate, so B1 checks the simulator's bookkeeping more than its predictive power. An independent verifier, a separate OpenAI Codex session that reused the prediction agents' copy of the repository rather than a fresh one, reran a subset, whose outputs matched exactly, and recomputed every B2 and B3 classification (Appendix~\ref{app:timeline}). A second round, B1n--B4n, repeated the four predictions with the native pool, ten seeds for B2n and B3n. B3n is blind in time; B2n was pushed before any of its runs had finished, but after their jobs were submitted; B1n and B4n are not blind, because the native curve and $\lmiss$ were already in the project status notes when the round was specified. No independent verifier reran this round (\S\ref{sec:blind}, Appendix~\ref{app:timeline}). A third round, made by the same kind of agent with the simulator unchanged and with its token costs scaled to vLLM's measured all-miss token rate (a fit, not a prediction), predicted vLLM's closed curve (B4v), from-empty edge (B2v) and paired test (B3v); it reached GitHub 22 minutes after the vLLM Stage~C and E0b jobs were submitted and before any of them had ended, and before the vLLM registration existed (\S\ref{sec:vllm}).

Our earlier simulator (rounds 1--5), which is not calibrated to the engine and uses assumed costs (all-miss capacity 0.0474\sps), is used for the mechanism comparison of \S\ref{sec:mech}, in units of its own $\lmiss$, and in \S\ref{sec:fluid}, \S\ref{sec:prereg} and \S\ref{sec:sim} where stated.
